%! Introduction to Polynomial Functions — Unit 2, Lesson 2.0 — Exit Ticket (Answer Key)
\documentclass[10pt]{article}
\usepackage{precalculus-article}
\usepackage{precalculus-key}   % pulls in -boxes; do NOT also load -boxes
\newcommand{\TallMath}[1]{$\displaystyle #1\rule[-1.4em]{0pt}{3.2em}$}

\begin{document}

\pageheader{Unit 2, Lesson 2.0}{Exit Ticket --- Answer Key}

\vspace{0.12in}

\begin{enumerate}[itemsep=12pt]

\item Consider the polynomial $p(x) = -4x^3 + 7x^2 - x + 9$.

\vspace{0.06in}
\begin{enumerate}[label=(\alph*), itemsep=6pt]
  \item Degree: \ans{3} \qquad Leading coefficient: \ans{$-4$}
        \qquad Constant term: \ans{9}
  \item By degree, $p$ belongs to the family called \ans{cubic}.
\end{enumerate}

\item Is each expression a polynomial? Circle one. Then explain the \textbf{No} in
      one sentence.

\begin{enumerate}[label=(\alph*), itemsep=6pt]
  \item $g(x) = 6x^4 - 2x + 1$ \quad \ans{Yes} / \textbf{No}
  \item $h(x) = 3x^2 + \dfrac{5}{x}$ \quad \textbf{Yes} / \ans{No}
\end{enumerate}

\vspace{0.06in}
Reason for the one that is \textbf{not} a polynomial:

\vspace{0.06in}
\ansline{In $h$ the variable sits in a denominator: $5/x = 5x^{-1}$, and a negative exponent is not allowed.}

\item Write $5 - 2x + 4x^2$ in standard form.

\vspace{0.1in}
Standard form: \ans{$4x^2 - 2x + 5$}

\end{enumerate}

\end{document}
